% !TEX program = xelatex
% !TEX options = -synctex=1 -interaction=nonstopmode -file-line-error
% example-scientific.tex
% Demonstrates rossbeamerscientific features: title page with \conference,
% \pval, \ci, finding/limitation boxes, study design, backup slides, \cite.
\documentclass[palette=lancet]{rossbeamerscientific}

\title{Sleep Architecture and Cognitive Decline}
\subtitle{A Prospective Cohort Study}
\author{R.~Dunne\textsuperscript{1}, J.~Smith\textsuperscript{2}}
\institute{\textsuperscript{1}University of Manchester \quad
           \textsuperscript{2}MRC Cognition and Brain Sciences Unit}
\date{March 2026}
\conference{International Neuropsychological Society, Vienna}

% Dummy bib for demonstration
\begin{filecontents}[overwrite]{\jobname.bib}
@article{smith2024,
  author  = {Smith, J. and Doe, A.},
  title   = {Slow-wave sleep and memory consolidation},
  journal = {Neurology},
  year    = {2024},
  volume  = {102},
  pages   = {1120--1130},
}
@article{jones2023,
  author  = {Jones, B. and Lee, C.},
  title   = {Age-related changes in sleep microstructure},
  journal = {Sleep},
  year    = {2023},
  volume  = {46},
  pages   = {zsad041},
}
\end{filecontents}
\addbibresource{\jobname.bib}

\begin{document}

% --- Title ---
\begin{frame}[plain]
  \titlepage
\end{frame}

% --- Research Question ---
\researchquestion{Does reduced slow-wave sleep predict domain-specific
  cognitive decline over a 3-year follow-up period?}

% --- Study Design ---
\begin{studydesign}{Study Design}
  \designitem{Design}{Prospective longitudinal cohort}
  \designitem{Sample}{$N = 412$ community-dwelling older adults, aged 60--85}
  \designitem{Exposure}{Polysomnography-derived slow-wave sleep (\% of TST)}
  \designitem{Outcome}{Change in RBANS domain scores at 36 months}
  \designitem{Analysis}{Linear mixed models adjusted for age, sex, education, and AHI}
\end{studydesign}

% --- Results ---
\begin{frame}{Primary Outcome}
  Participants in the lowest SWS tertile showed significantly greater decline
  in Delayed Memory compared to the highest tertile.\cite{smith2024}

  \vskip8pt
  \begin{itemize}
    \item Delayed Memory: \ci{$\beta = -4.2$}{$-6.8$}{$-1.6$}, \pval{0.002}
    \item Immediate Memory: \ci{$\beta = -2.1$}{$-4.5$}{$0.3$}, \pval{0.09}
    \item Visuospatial: \ci{$\beta = -0.8$}{$-3.0$}{$1.4$}, \pval{0.48}
  \end{itemize}
  \sigmarkers
\end{frame}

% --- Finding box ---
\begin{frame}{Key Findings}
  \begin{findingbox}
    Reduced slow-wave sleep was selectively associated with accelerated
    decline in delayed memory, but not visuospatial or language domains,
    after controlling for sleep-disordered breathing.\cite{jones2023}
  \end{findingbox}

  \vskip6pt
  \begin{limitationbox}
    Single-night polysomnography may not capture habitual sleep patterns.
    Replication with multi-night actigraphy is warranted.
  \end{limitationbox}

  \vskip6pt
  \begin{implicationbox}
    Sleep interventions targeting SWS enhancement may offer a modifiable
    pathway to delay memory decline in ageing populations.
  \end{implicationbox}
\end{frame}

% --- Equation example ---
\begin{frame}{Statistical Model}
  The primary model was specified as:
  \[
    Y_{ij} = \beta_0 + \beta_1 \cdot \text{SWS}_i
             + \eqhl{\beta_2 \cdot \text{Time}_j}
             + \beta_3 \cdot (\text{SWS}_i \times \text{Time}_j)
             + \mathbf{X}_i \boldsymbol{\gamma} + u_i + \varepsilon_{ij}
  \]
  \slidefootnote{Model estimated using restricted maximum likelihood (REML).}
\end{frame}

% --- References ---
\begin{referenceslide}
  \printbibliography[heading=none]
\end{referenceslide}

% --- Backup ---
\backupslides

\begin{frame}{Sensitivity Analysis}
  Excluding participants with AHI~$>$~15 did not materially change results:
  \begin{itemize}
    \item Delayed Memory: \ci{$\beta = -3.8$}{$-6.5$}{$-1.1$}, \pval{0.006}
  \end{itemize}
\end{frame}

\end{document}
